% =============================================================================
% Title Block (spans both columns via \twocolumn[...])
% =============================================================================

\twocolumn[{
\begin{flushleft}

{\small
Department of Computer Engineering\\
Faculty of Sciences and Techniques Tangier\\
Abdelmalek Essaadi University, Morocco\\
2025/2026\par
}

\vspace{2.5em}

{\color{titleblue}\Huge\bfseries
Système de Question-Réponse Juridique Marocain\\[0.2em]
Basé sur le Retrieval-Augmented Generation (RAG)\\[0.2em]
avec Traitement du Langage Naturel Arabe\par
}

\vspace{1.5em}

{\large\bfseries
Encadré par : Pr. BENABDELOUAHAB Ikram\\[0.8em]
Réalisé par : ZARQI Ezzoubair, EL ATTAR Mohamed, EL GHARIB Mahmoud, HAMDAN
Kaouthar ET AMAMI Yousra\par
}


\end{flushleft}

\vspace{1.5em}

% Abstract
\begin{leftbar}
\textbf{\textcolor{headingblue}{RÉSUMÉ}} Cet article présente la conception,
l'implémentation et l'évaluation d'un système de Question-Réponse (Q\&R)
juridique spécialisé dans le droit marocain, basé sur l'architecture
Retrieval-Augmented Generation (RAG). Le système intègre une chaîne complète
de traitement du langage naturel arabe incluant l'extraction de texte à partir
de documents PDF (numériques et scannés), le nettoyage, la normalisation
arabe, l'analyse structurelle des articles de loi, et l'indexation vectorielle
via PostgreSQL avec l'extension pgvector. La recherche hybride combine la
similarité cosinus dense (modèle \texttt{BAAI/bge-m3}, 1024 dimensions) avec
la recherche lexicale BM25, fusionnée par Reciprocal Rank Fusion (RRF).
Le module de génération s'appuie sur le modèle Gemma 4 2B (Google), servi
sur GPU via Google Colab et accessible par tunnel Cloudflare. Les benchmarks
montrent que la stratégie RRF atteint un Hit@5 de 68,75\% et un Recall@5 de
59,38\%, tandis que la recherche dense offre le meilleur MRR@5 (0,565).
L'évaluation des modèles d'embedding démontre la supériorité des modèles
1024 dimensions (\texttt{intfloat/multilingual-e5-large}, \texttt{BAAI/bge-m3})
pour le domaine juridique arabe, avec une précision de domaine atteignant
89\%. Ce travail constitue la première plateforme open-source complète de
RAG juridique marocain, couvrant six domaines de droit (pénal, trafic,
familial, civil, commercial, procédures).
\end{leftbar}

\vspace{1em}

% Keywords
\begin{leftbar}
\textbf{\textcolor{headingblue}{MOTS-CLÉS}} RAG, NLP Arabe, Droit Marocain,
Recherche Hybride, Embeddings Multilingues, pgvector, Gemma 4,
Question-Réponse Juridique, Reciprocal Rank Fusion.
\end{leftbar}

\vspace{2.5em}

}]